\documentclass[11pt]{article}
\usepackage[margin=1in]{geometry}
\usepackage{booktabs}
\title{DME: SNc vs.\ VTA, within Healthy Baseline --- Interpretation}
\author{GSE233866, untreated cohort}
\date{}

\begin{document}
\maketitle

\subsection*{Why this closes the loop}

\texttt{snc\_vs\_vta\_preservation/} (same folder) already established that SNc's
CACNA1D-containing module structurally reproduces in VTA (Zsummary=20.4) --- a test of
whether the same \emph{gene community} exists in both regions. It does not test whether
that module's \emph{activity level} differs between regions, because \texttt{ModulePreservation}
and \texttt{FindDMEs} are different statistics entirely. This network runs the missing test:
pool healthy-baseline SNc and VTA cells (no lesion arm, no disease model, the cleanest
possible comparison this project can make) into one network and test whether CACNA1D's
module differs in level between the two regions.

\subsection*{Result: the strongest signal in the whole project}

CACNA1D's module here is \textbf{yellow}, and it is differentially active between SNc and
VTA with $\mathrm{avg\_log2FC}=+14.55$, $p_{\mathrm{adj}}=4.17\times10^{-212}$, higher in
SNc --- the single most significant module in this 11-module table, more extreme than any
other module tested here, and comparable in magnitude to (in fact slightly larger than) the
equivalent results from the intact (+13.21) and lesioned (+13.34) states tested elsewhere in
this project.

\subsection*{What this settles}

Combined with \texttt{../../lesion\_model/snc\_vs\_vta\_combined\_dme\_intact/} and
\texttt{snc\_vs\_vta\_combined\_dme\_lesioned/}, this is now a three-for-three result: CACNA1D's
module is significantly more active in SNc than VTA in the healthy state, the intact
lesion-arm state, and the lesioned state alike, at comparable magnitude in all three. This
closes the question raised earlier in this project's own history --- whether the SNc-vs-VTA
gap seen during lesioning was disease-induced. It is not. It is present with \emph{zero}
disease model involved, in a completely separate cohort of animals from the lesion arm. The
correct framing for this project's central region-selectivity claim is now: \textbf{SNc and
VTA are constitutively distinct in CACNA1D's co-expression program, independent of disease
state}, which the lesion-arm data shows persists (and is not obviously amplified) through
6-OHDA lesioning.

\subsection*{Caveats}

As with every DME table in this project, treat log2FC magnitude as directional evidence, not
a literal fold-change --- module eigengenes can be negative or near-zero, which destabilizes
that arithmetic independent of whether the $p$-value is sound. This network pools only 3,058
cells (2,226 SNc + 832 VTA), the smallest of the three SNc-vs-VTA DME networks in this
project, reflecting the healthy-baseline cohort's smaller DA-neuron yield rather than a
weaker analysis --- the $p$-value here is nonetheless the most extreme of the three, so this
is not a power limitation in practice.

\end{document}
